% ---------------------------------------------------------------------------
% SICEDU - Guia de pruebas manuales
% Compilar con: pdflatex GUIA_PRUEBAS.tex   (dos veces, por el indice)
% ---------------------------------------------------------------------------
\documentclass[11pt,a4paper]{article}

\usepackage[utf8]{inputenc}
\usepackage[T1]{fontenc}
\usepackage[spanish,es-noquoting,es-tabla]{babel}
\usepackage[margin=2.3cm,headheight=14pt]{geometry}
\usepackage{lmodern}
\usepackage{amssymb}
\usepackage{booktabs}
\usepackage{tabularx}
\usepackage{longtable}
\usepackage{enumitem}
\usepackage{xcolor}
\usepackage{colortbl}
\usepackage{fancyhdr}
\usepackage{titlesec}

% Los colores se definen ANTES de hyperref: los recibe como opciones y tienen
% que existir ya cuando se carga.
\definecolor{azulmarino}{HTML}{14274E}
\definecolor{grisclaro}{HTML}{F4F6F9}
\definecolor{ambar}{HTML}{B45309}

\usepackage[colorlinks=true,linkcolor=azulmarino,urlcolor=azulmarino]{hyperref}

\titleformat{\section}{\Large\bfseries\color{azulmarino}}{\thesection.}{0.6em}{}
\titleformat{\subsection}{\large\bfseries}{\thesubsection}{0.6em}{}

\pagestyle{fancy}
\fancyhf{}
\lhead{\color{azulmarino}\textbf{SICEDU} \textperiodcentered\ Guia de pruebas}
\rhead{\thepage}
\renewcommand{\headrulewidth}{0.4pt}

\newcommand{\paso}[1]{\textbf{#1}}
\newcommand{\ruta}[1]{\textcolor{azulmarino}{\texttt{#1}}}
\newcommand{\ok}{$\square$}
\newcommand{\flecha}{$\rightarrow$}

% Aviso destacado.
%
% Va con minipage dentro de una caja, no con tabularx: en una celda de párrafo
% el \\ final rompe la linea DENTRO de la celda en vez de cerrar la fila, y el
% \hline siguiente falla. Asi admite listas y parrafos sin sorpresas.
\newsavebox{\cajaaviso}
\newenvironment{aviso}[1]
  {\par\medskip\noindent
   \begin{lrbox}{\cajaaviso}
   \begin{minipage}{\dimexpr\textwidth-2\fboxsep-2\fboxrule\relax}
   \smallskip
   \textbf{\textcolor{ambar}{#1}}\par\smallskip}
  {\smallskip\end{minipage}\end{lrbox}%
   \fcolorbox{ambar}{grisclaro}{\usebox{\cajaaviso}}\par\medskip}

\setlist[itemize]{itemsep=2pt,topsep=4pt}
\setlist[enumerate]{itemsep=2pt,topsep=4pt}

\begin{document}

% ---------------------------------------------------------------------------
\begin{titlepage}
\centering
\vspace*{3cm}
{\Huge\bfseries\color{azulmarino} SICEDU\par}
\vspace{0.4cm}
{\LARGE Guía de pruebas manuales\par}
\vspace{0.3cm}
{\large Sprint de cierre: seguridad, offline y mantenimientos\par}
\vspace{2cm}
{\large Misión Huascarán \textperiodcentered\ Frontend\par}
\vspace{0.3cm}
{\large 7 de octubre de 2026\par}
\vfill
{\normalsize Dirigido a quienes ejecutan las pruebas.\par}
{\normalsize No hace falta saber programar para seguirla.\par}
\end{titlepage}

\tableofcontents
\newpage

% ---------------------------------------------------------------------------
\section{Antes de empezar}

Esta guía recorre lo construido en el sprint de cierre. Cada apartado dice
\emph{qué hacer} y \emph{qué debe ocurrir}. Si lo que ocurre no coincide,
anótelo en la hoja de incidencias del final (página~\pageref{sec:incidencias}).

\subsection{Dos formas de levantar la aplicación}

\paragraph{Opción A --- sin backend.} Es la más cómoda para casi todo. En el
fichero \texttt{.env} deje:

\begin{quote}\ttfamily
VITE\_AUTH\_REAL=false\\
VITE\_ADMIN\_REAL=false
\end{quote}

Reinicie con \texttt{npm run dev}. La contraseña de todas las cuentas es
\texttt{sicedu123}:

\begin{center}
\begin{tabular}{ll}
\toprule
\textbf{Rol} & \textbf{Correo} \\
\midrule
Docente    & \texttt{rcardenas@sicedu.test}  \\
Supervisor & \texttt{jefatura@sicedu.test}   \\
Directivo  & \texttt{direccion@sicedu.test}  \\
\bottomrule
\end{tabular}
\end{center}

\paragraph{Opción B --- con backend.} Necesaria solo para lo que toca permisos
reales y para la expiración de sesión (apartado~\ref{sec:cu007}). Deje los dos
flags en \texttt{true}, levante el backend y use las cuentas reales
(\texttt{profesor.prueba@sicedu.test} / \texttt{ProfesorTest123}).

\subsection{Cómo simular que no hay conexión}

En el navegador: \textbf{F12} \flecha{} pestaña \textbf{Network} \flecha{} el
desplegable de \emph{throttling} \flecha{} \textbf{Offline}.

Para mirar lo que la aplicación guarda en el dispositivo:
\textbf{F12} \flecha{} \textbf{Application} \flecha{} \textbf{IndexedDB}.

\begin{aviso}{El orden importa en todo lo que sea sin conexión}
Hay que pulsar \emph{Iniciar actividad} \textbf{con conexión}, esperar a que
termine la descarga, y \emph{después} ponerse en modo Offline.

Si se inicia la actividad ya sin red no hay nada que descargar, y la aplicación
lo avisa con el mensaje «Sin conexión al iniciar». En ese caso las pantallas
saldrán vacías, y \textbf{es el comportamiento correcto}, no un fallo.
\end{aviso}

% ---------------------------------------------------------------------------
\newpage
\section{T15 --- Permisos y manejo de errores}

\subsection{El menú cambia según quién entra}

Entre con los tres roles y compare el menú lateral:

\begin{center}
\begin{tabularx}{\textwidth}{lX}
\toprule
\textbf{Rol} & \textbf{Debe ver} \\
\midrule
Docente    & Inicio, Alumnos, Asistencia, Rúbrica, Seguimiento de Lectura, Registro de Vuelo, Sesiones \\
Supervisor & Inicio, Usuarios, Colegios, Asistencia, Rúbrica, Seguimiento de Lectura, Registro de Vuelo, Sesiones, Dashboard \\
Directivo  & Solo Inicio y Dashboard \\
\bottomrule
\end{tabularx}
\end{center}

\paragraph{Ojo con Alumnos.} El Supervisor \emph{no} tiene «Alumnos» como
entrada propia del menú: la ve dentro de \ruta{Usuarios}, y solo de lectura. Eso
es intencionado, no un olvido.

\subsection{Los botones que no tocan, no están}

Entre como \textbf{Supervisor} en \ruta{Rúbrica}. No debe haber ningún control
editable ni botón «Crear grilla». El Supervisor mira, no escribe.

\subsection{Entrar sin internet}

Póngase en modo Offline \emph{antes} de entrar y recargue \ruta{/login}.

\textbf{Esperado:} el mensaje «Se requiere conexión para ingresar» y el botón de
entrar deshabilitado.

\subsection{Sesión caducada (error 401)}

\begin{enumerate}
  \item Escriba algo en una grilla.
  \item \textbf{F12} \flecha{} \textbf{Application} \flecha{} \textbf{Session
        Storage}: borre el token.
  \item Haga cualquier acción en la aplicación.
\end{enumerate}

\textbf{Esperado:} vuelve al login. Y, muy importante, en \textbf{IndexedDB}
\flecha{} \texttt{sicedu.cola-envios} \textbf{los cambios siguen ahí}. Perder la
sesión no puede perder el trabajo.

\subsection{Sobre el error 403}

Conviene saberlo para no perder tiempo buscándolo: es difícil de provocar
haciendo clic. Como el menú y los botones salen de la misma matriz de permisos,
la interfaz ya no ofrece ningún camino que acabe en 403. El manejo que hay es
una red de seguridad para una petición manual, no un recorrido de usuario.

% ---------------------------------------------------------------------------
\newpage
\section{T16 y T17 --- Mantenimientos}

Entre como \textbf{Supervisor} en \ruta{Usuarios}. Repita después lo mismo en
\ruta{Colegios} y en la pestaña \ruta{Alumnos}: las tres pantallas se comportan
igual a propósito.

\begin{enumerate}
  \item \paso{Estado = Activo por defecto.} El filtro arranca ahí sin tocarlo.
  \item \paso{Nuevo.} Abre un formulario vacío. Pulse \emph{Guardar} sin
        rellenar nada: debe salir «Este campo es obligatorio» debajo de
        \textbf{cada} campo, no un mensaje genérico arriba.
  \item \paso{Editar.} El mismo formulario, con los datos cargados. El campo
        \emph{Rol} aparece bloqueado.
  \item \paso{Ver.} El mismo formulario con todo bloqueado y \emph{sin} botón de
        guardar.
  \item \paso{Inactivar.} Pide confirmación y avisa de que no se borra nada.
        \begin{itemize}
          \item Si cancela: no pasa nada.
          \item Si confirma: la fila desaparece con el filtro \emph{Activos} y
                \textbf{reaparece} eligiendo \emph{Inactivos}.
        \end{itemize}
  \item \paso{Auditoría.} Dentro de \emph{Editar} o \emph{Ver}, pestaña
        \emph{Auditoría}: quién, cuándo, qué campo, valor antes y valor después.
\end{enumerate}

\subsection{Lo que añade T17}

\begin{center}
\begin{tabularx}{\textwidth}{lX}
\toprule
\textbf{Pantalla} & \textbf{Qué comprobar} \\
\midrule
Alumnos   & En el formulario hay \emph{Ciclo evaluado}, separado de \emph{Grado}. Un alumno de 6.\textordmasculine{} puede evaluarse con la rúbrica del ciclo~III. \\
\addlinespace
Usuarios  & Al elegir rol \emph{Docente} aparecen \emph{Colegios asignados} y \emph{Secciones asignadas}, con cada sección etiquetada con su colegio. \\
\addlinespace
Colegios  & \emph{Grados y secciones} ya no es un cuadro de texto libre: es una lista con botón \emph{Añadir sección} y un grado por cada una. \\
\bottomrule
\end{tabularx}
\end{center}

% ---------------------------------------------------------------------------
\newpage
\section{T18 y T19 --- Actividad de trabajo y modo sin conexión}
\label{sec:offline}

El botón de actividad está \textbf{arriba a la derecha} y es solo para el
Docente. Con Supervisor o Directivo no debe aparecer.

\subsection{Sin actividad abierta no se puede escribir}

Como Docente, sin pulsar nada, entre en \ruta{Rúbrica}.

\textbf{Esperado:} el aviso «Pulse Iniciar actividad\ldots» y los selectores
deshabilitados. Al pulsar el botón, se habilitan.

\subsection{Las horas de la sesión}

Tras iniciar, vaya a \ruta{Inicio}. Debe leerse algo como «Sesión iniciada a las
08:15. Finaliza a las 16:15.»

\textbf{Recargue la página:} las horas \textbf{no} deben cambiar. Si cambian, la
sesión se está reiniciando sola, que es justo lo que no debe pasar.

\subsection{La descarga previa}

Pulse \emph{Iniciar actividad} \textbf{con conexión}. Después, en
\textbf{IndexedDB}, busque \texttt{sicedu.precarga}.

\textbf{Esperado:} debe contener \texttt{alumnos}, \texttt{colegios},
\texttt{grados}, \texttt{semanas}, \texttt{asignaciones} y los catálogos de
niveles.

Si sale el aviso «Precarga incompleta», algo no se pudo descargar: anótelo,
porque es exactamente lo que dejará huecos en los pasos siguientes.

\subsection{Trabajar sin conexión}

\begin{enumerate}
  \item Con la actividad ya iniciada y descargada, póngase en \textbf{Offline}.
  \item Arriba aparece la etiqueta «Sin conexión».
  \item Entre en \ruta{Rúbrica}. \textbf{Los desplegables de colegio y semana
        deben tener sus opciones}, sacadas de lo descargado antes.
  \item Edite una fila. La etiqueta pasa a «Sin conexión \textperiodcentered{} 1
        pendiente».
\end{enumerate}

\begin{aviso}{Esto es lo que fallaba hasta ahora}
El punto 3 es nuevo. Antes, al entrar en Rúbrica sin conexión, el desplegable de
colegios salía vacío y no se podía abrir ninguna grilla, aunque los colegios
estuvieran descargados. Merece la pena probarlo con calma.
\end{aviso}

\subsection{Avisos al salir}

Con cambios pendientes, pulse su nombre (abajo a la izquierda) \flecha{}
\emph{Finalizar sesión}.

\textbf{Esperado:} un aviso de que no se va a perder nada.

\subsection{La limpieza condicional}

Esto comprueba una regla delicada: los datos de menores se borran del
dispositivo, \textbf{pero nunca si hay trabajo sin enviar}.

\begin{enumerate}
  \item Salga \textbf{con} cambios pendientes \flecha{} en IndexedDB,
        \texttt{sicedu.cola-envios} \textbf{sigue existiendo}.
  \item Vuelva a conectar, espere a que la etiqueta diga «Todo sincronizado», y
        salga otra vez \flecha{} ahora sí \textbf{desaparece}.
\end{enumerate}

% ---------------------------------------------------------------------------
\newpage
\section{T20 --- Sesiones}

Como \textbf{Docente} solo se ven las propias. Como \textbf{Supervisor} aparecen
tres filtros --- docente, colegio y fecha --- y un botón \emph{Ver detalle} en
cada fila.

En el detalle debe verse, por cada cambio: el registro afectado, el campo, el
valor anterior y el nuevo.

\paragraph{Lo interesante.} Busque la etiqueta naranja \emph{«Sincronizado de
forma diferida»}. Aparece solo en los cambios cuya hora de captura es distinta
de la de recepción, es decir, los que se hicieron sin conexión y llegaron
después. Es la forma de demostrar que el trabajo offline queda trazado.

% ---------------------------------------------------------------------------
\section{T22, T30 y T32 --- Grillas del aula}

\subsection{Asistencia (T30)}

Entre como Docente en \ruta{Asistencia}. Los filtros van en este orden:
colegio \flecha{} ciclo \flecha{} grado \flecha{} sección \flecha{} fecha.

\begin{itemize}
  \item \textbf{Al entrar, el colegio y el grado ya vienen elegidos}, y el
        desplegable de \emph{Sección} trae las secciones de ese colegio. No debe
        quedarse ninguno en blanco.
  \item Ponga una \textbf{fecha futura}: sale «No existe una grilla registrada
        para esta fecha» con el botón para crearla.
  \item Créela: debe nacer con una fila por alumno y \textbf{vacía}. No puede
        copiar los datos de otro día.
  \item Marque \emph{Presente} / \emph{Ausente} y guarde.
\end{itemize}

\begin{aviso}{Esto también acaba de arreglarse}
Antes, al entrar en Asistencia, la pantalla \emph{aparentaba} tener un colegio
elegido pero no lo tenía: el desplegable de Sección salía vacío y el cuerpo
seguía pidiendo «Elija colegio, grado y fecha». Compruebe que ya no pasa.
\end{aviso}

\subsection{Rúbrica (T32)}

Mismo orden de filtros. Elija una semana sin datos y debe salir «Crear grilla
para esta semana».

\subsection{La ausencia enlazada}

Esta es la comprobación con más valor del apartado, porque cruza dos módulos:

\begin{enumerate}
  \item En \ruta{Asistencia}, marque a un alumno como \textbf{Ausente} y guarde.
  \item Vaya a \ruta{Rúbrica}, misma semana y mismo grado.
  \item Ese alumno debe aparecer con la etiqueta roja \textbf{Ausente}, y sus
        selectores de \emph{Fluidez} y \emph{Comprensión} deben estar
        \textbf{deshabilitados}.
\end{enumerate}

La idea es que la falta se marque \textbf{una sola vez}, en Asistencia, y que
nadie pueda evaluar a quien no estuvo.

\subsection{Solo lectura del Supervisor (T22)}

Entre como Supervisor en \ruta{Asistencia} o \ruta{Rúbrica}: aparece el filtro
\emph{Docente}, y no hay edición ni botón de crear grilla.

% ---------------------------------------------------------------------------
\newpage
\section{CU007 --- Expiración a las 8 horas}
\label{sec:cu007}

\begin{aviso}{Esto no se puede probar sin backend}
En modo simulado el token es \texttt{mock.4.2026}, que no es un JWT real y no
lleva fecha de caducidad, así que el temporizador nunca llega a armarse. No es
un fallo: es que no hay nada que medir.
\end{aviso}

Para probarlo de verdad, en el \texttt{.env} \textbf{del backend}:

\begin{quote}\ttfamily
ACCESS\_TOKEN\_EXPIRE\_MINUTES=1
\end{quote}

Reinicie el backend, entre con \texttt{VITE\_AUTH\_REAL=true} y espere un minuto
sin tocar nada.

\textbf{Esperado:} la aplicación lleva al login con el aviso «Su sesión terminó
por alcanzar su tiempo máximo de 8 horas».

\textbf{Acuérdese de dejarlo otra vez en 480} al terminar.

% ---------------------------------------------------------------------------
\section{Lo que verán simulado}

Conviene saberlo de antemano para no reportar como fallo algo que todavía no
existe. Estas pantallas funcionan y se ven completas, pero sus datos \textbf{no
vienen del backend}, porque esos endpoints aún no están publicados:

\begin{itemize}
  \item Auditoría
  \item Sesiones
  \item Secciones
  \item Asistencia
  \item Los resúmenes del módulo de Inicio
\end{itemize}

Cada una se conecta cambiando una línea cuando el backend las publique. Lo que
sí tiene sentido probar en ellas es el \emph{comportamiento}: los filtros, los
estados vacíos, los permisos y los avisos.

% ---------------------------------------------------------------------------
\newpage
\section{Hoja de incidencias}
\label{sec:incidencias}

Una fila por cada cosa que no coincida con lo esperado. Cuanto más concreto el
«qué ocurrió», más rápido se arregla.

% longtable se centra sola; meterla en un center descuadra el espaciado.
\begin{longtable}{p{0.8cm}p{3.2cm}p{4.8cm}p{4.8cm}}
\toprule
\textbf{\ok} & \textbf{Apartado} & \textbf{Qué debía ocurrir} & \textbf{Qué ocurrió} \\
\midrule
\endhead
\ok & & & \\[0.8cm]
\ok & & & \\[0.8cm]
\ok & & & \\[0.8cm]
\ok & & & \\[0.8cm]
\ok & & & \\[0.8cm]
\ok & & & \\[0.8cm]
\ok & & & \\[0.8cm]
\ok & & & \\[0.8cm]
\ok & & & \\[0.8cm]
\ok & & & \\[0.8cm]
\bottomrule
\end{longtable}

\vfill
\noindent\rule{\textwidth}{0.4pt}

\noindent\small Al reportar, indique siempre tres cosas: el \textbf{rol} con el
que entró, si estaba en \textbf{modo simulado o con backend}, y si tenía o no
\textbf{conexión}. Las tres cambian el comportamiento esperado.

\end{document}
